\chapter{耶稣的神迹—— (David/YingYing)}

\section{医治四人抬的瘫子}
\subsection{中文}

\subsubsection{創世記 1:3,6-7,9}
神說：「要有光。」就有了光。\\
神說：「諸水之間要有空氣，將水分為上下。」 7 神就造出空氣\\
神說：「天下的水要聚在一處，使旱地露出來。」事就這樣成了。
\subsubsection{希伯来书 9:22}
按着律法，几乎每样东西都是用血洁净的；没有流血，就没有赦罪。
\subsubsection{约壹1:7}
「祂儿子耶稣的血，也洗净我们一切的罪。」
\subsubsection{可 2:1-12}
過了些日子，耶穌又進了迦百農。人聽見他在房子裡， 2 就有許多人聚集，甚至連門前都沒有空地。耶穌就對他們講道。 3 有人帶著一個癱子來見耶穌，是用四個人抬來的。 4 因為人多，不得近前，就把耶穌所在的房子拆了房頂。既拆通了，就把癱子連所躺臥的褥子都縋下來。 5 耶穌見他們的信心，就對癱子說：「小子，你的罪赦了！」 6 有幾個文士坐在那裡，心裡議論說： 7 「這個人為什麼這樣說呢？他說僭妄的話了！除了神以外，誰能赦罪呢？」 8 耶穌心中知道他們心裡這樣議論，就說：「你們心裡為什麼這樣議論呢？ 9 或對癱子說『你的罪赦了』，或說『起來，拿你的褥子行走』，哪一樣容易呢？ 10 但要叫你們知道，人子在地上有赦罪的權柄。」就對癱子說： 11 「我吩咐你，起來，拿你的褥子回家去吧！」 12 那人就起來，立刻拿著褥子，當眾人面前出去了。以致眾人都驚奇，歸榮耀於神，說：「我們從來沒有見過這樣的事！」


\newpage
\Large
\subsection{\LARGE{Teacher's Version - English Version}}
\large
\vspace{.25cm}
\subsubsection{\large{Genesis 1:3,6-7,9/Hebrews 9:22/1 John 1:7 (ICB)}}
【Genesis 1:3,6-7,9 (ICB)】God said, “Let there be light!” And there was light. ...Then God said, “Let there be something to divide the water in two!” So God made the air to divide the water in two. ...Then God said, “Let the water under the sky be gathered together so the dry land will appear.” And it happened.\\
【Hebrews 9:22 (ICB)】 The law says that almost everything must be made clean by blood. And sins cannot be forgiven without blood to show death.\\
【1 John 1:5,7 (ICB)】 God is light, and in him there is no darkness at all. And when we live in the light, the blood of the death of Jesus, God’s Son, is making us clean from every sin.\\
\vspace{.25cm}
\textbf{Questions}\vspace{-.5cm}
\begin{itemize}[noitemsep]
\item According to Genesis, how did God create the universe? \textbf{\underline{by word/speaking}}
\item According to the Law, what should be used to make everything clean? \textbf{\underline{by blood}}
\item How can sins be forgiven? \textbf{\underline{blood to show death}}
\item Whose blood of the death can make us clean from every sin? \textbf{\underline{the death of Jesus, God's Son}}
\item What is light according to John? \textbf{\underline{God is light}}
\item What should we do to be cleaned? \textbf{\underline{when we live in the light}}
\end{itemize}

\vspace{.25cm}
\subsubsection{\large{Mark 2:1-3}}
A few days later, Jesus came back to Capernaum. The news spread that he was home. So many people gathered to hear him preach that the house was full. There was no place to stand, not even outside the door. Jesus was teaching them. Some people came, bringing a paralyzed man to Jesus. Four of them were carrying the paralyzed man. \\
\vspace{.25cm}
\textbf{Questions}\vspace{-.5cm}
\begin{itemize}[noitemsep]
\item “he was home”， here, "home" refers to which city? \textbf{\underline{Capernaum}}
\item Where did Jesus give his preach? \textbf{\underline{in the houe}}
\item Why there was no place to stand? \textbf{\underline{So many people gathered in the house}}
\item Why the man was carried by others? \textbf{\underline{paralyzed man}}
\item How many people carried the man? \textbf{\underline{4}}
\end{itemize}

\vspace{.25cm}
\subsubsection{\large{Mark 2:4-5}}
But they could not get to Jesus because of the crowd. So they went to the roof above Jesus and made a hole in the roof. Then they lowered the mat with the paralyzed man on it. Jesus saw that these men had great faith. So he said to the paralyzed man, “Young man, your sins are forgiven.” \\
\vspace{.25cm}
\textbf{Questions}\vspace{-.5cm}
\begin{itemize}[noitemsep]
\item How did the four man send the paralyzed man to Jesus?\\  \textbf{\underline{made a hole in the roof / lowered him with the mat}}
\item Whose great faith was found by Jesus? From which description you draw the conclusion?  \textbf{\underline{4 men and the paralyzed man }}
\item The paralyzed man was a \textbf{\underline{ young man}}\\
A, old man \hspace{1cm} B, child, \hspace{1cm} C, young man
\item What did Jesus say to the paralyzed man? \textbf{\underline{Young man, your sins are forgiven.}}
\end{itemize}

\newpage
\vspace{.3cm}
\subsubsection{\large{Mark 2:6-7}}
Some of the teachers of the law were sitting there. They saw what Jesus did, and they said to themselves, “Why does this man say things like that? He is saying things that are against God. Only God can forgive sins.” \\
\vspace{.2cm}
\textbf{Questions}\vspace{-.5cm}
\begin{itemize}[noitemsep]
\item Who were sitting with Jesus in the room, and said to themselves? \textbf{\underline{ teachers of the law}}
\item“Why does this man say things like that"? here, "\textbf{this man}“ refers to whom? \textbf{\underline{ Jesus}}
\item“Why does this man say things like that"? here, "\textbf{things}“ refers to what? \textbf{\underline{ "your sins are forgiven"}}
\item Why did they disagree with Jesus? \textbf{\underline{Only God can forgive sins, while man cannot}}
\item When these teachers thought "Only God can forgive sins", who did they think Jesus is?\textbf{\underline{a man}}
\end{itemize}


\vspace{.3cm}
\subsubsection{\large{Mark 2:8-12}}
At once Jesus knew what these teachers of the law were thinking. So he said to them, “Why are you thinking these things? Which is easier: to tell this paralyzed man, ‘Your sins are forgiven,’ or to tell him, ‘Stand up. Take your mat and walk’? But I will prove to you that the Son of Man has authority on earth to forgive sins.” So Jesus said to the paralyzed man, “I tell you, stand up. Take your mat and go home.” Immediately the paralyzed man stood up. He took his mat and walked out while everyone was watching him. The people were amazed and praised God. They said, “We have never seen anything like this!”\\ 
\vspace{.25cm}
\textbf{Questions}\vspace{-.5cm}
\begin{itemize}[noitemsep]
\item Did Jesus know what's the thought in these teachers's mind?\underline{yes}\\

\item "I will prove to you that the Son of Man has authority on earth", here \textbf{the Son of Man} refers to whom? The word of \textbf{authority} refers to which kind of authority on earth?\underline{Jesus/forgive sins}\\

\item The paralyzed man was brought to Jesus, what was this man's epectation from Jesus? \textbf{\underline{healing}} Did Jesus know this man's expectation?\underline{yes}\\

\item However, what \textbf{did} Jesus want to give the paralyzed man prior to giving him the healling?\\
\underline{forgive his sins}\\

\item Why did Jesus say： “Young man, your sins are forgiven” rather than say ‘Stand up. Take your mat and walk’ first? (Use the word of Bible to answer the question)\\
\underline{I will prove to you that the Son of Man has authority on earth to forgive sins.}\\

\item Jesus asked a question to people in the house, but Jesus did not give the answer. What is your answer about Jesus's question. 【Which is easier: to tell this paralyzed man, ‘\textbf{Your sins are forgiven},’ or to tell him, ‘\textbf{Stand up. Take your mat and walk}’?】\\
\underline{这里学生可以自由发挥，不用纠正，后面有更详细的引导讨论题。Student can speak freely}\\

\item If \underline{a man} tells the paralyzed man, ‘\textbf{Stand up. Take your mat and walk}’, how can others determin his statement has the power of healing or not?\\
\underline{By watching if the paralyzed man can stand up or not !!! }\\

\item If \underline{a man} tells the paralyzed man, ‘\textbf{your sins are forgiven}’, can others easily determin his word has the power of forgiven or not? or his statement is "TRUE" or "FALSE"?\\
\underline{No, thre is no way to tell if the man's word is 'TRUE' or 'FALSE'.}\\

\item For Jesus, what he has to do to make ‘\textbf{your sins are forgiven}’ become TRUE?\\
\underline{Die on the cross'.}\\

\item What is the price Jesus has to pay to get the power of "forgiven sins"?
\underline{Die on the cross'.}\\


\item How did Jesus prove to the people on site that he really has the power of  "forgiven sins"? or why did the sequence of Jesus' word could prove that he really has the power of  "forgiven sins"?\\
%\textbf{\underline{ANSWER}}\\
\underline{By saying ‘\textbf{Your sins are forgiven},’ and then say  ‘\textbf{Stand up. Take your mat and walk}’} \\
\underline{The logic behand here was: if Jesus was NOT God, but pretended to be God by saying}\\
\underline{‘\textbf{Your sins are forgiven},’ then Jesus 2nd word ‘\textbf{Stand up. Take your mat and walk}’ }\\
\underline{would not work, since God does not listen to sinners' prayer!!!. }\\
\underline{However, if Jesus was the TRUE God, the paralyzed man stood up would prove that Jesus' }\\
\underline{first statement of ‘\textbf{Your sins are forgiven},’ was TRUE!!!}\\


\item How did God create the universe? What did Jesus do to heal the paralyzed man?\\
\textbf{\underline{by word/by speaking}}\\

\item \textbf{For Jesus} which is easier: to tell this paralyzed man, ‘\textbf{Your sins are forgiven},’ or to tell him, ‘\textbf{Stand up. Take your mat and walk}’?\\
\underline{‘\textbf{Your sins are forgiven},’ is harder, since he has to die on the cross!!!}\\
\underline{‘\textbf{Stand up. Take your mat and walk}’ is much easier, just by speaking}\\


\item \textbf{For human being} which is easier: to tell this paralyzed man, ‘\textbf{Your sins are forgiven},’ or to tell him, ‘\textbf{Stand up. Take your mat and walk}’?\\
\underline{‘\textbf{Your sins are forgiven},’ is easier, since there is no way to prove this word is TRUE or FALSE!!!}\\
\underline{‘\textbf{Stand up. Take your mat and walk}’ is much hearder, since no one can heal the paralyzed}\\
\underline{man by speaking}\\

\item What were the people's reactions after seeing these?\\
\underline{The people were amazed and praised God}

\item When our prayer is different from God's desire, what do we learn from this passage?
\end{itemize}












\newpage
\Large
\subsection{\LARGE{Sudent Version - English Version}}
\large
\vspace{.25cm}
\subsubsection{\large{Genesis 1:3,6-7,9/Hebrews 9:22/1 John 1:7 (ICB)}}
【Genesis 1:3,6-7,9 (ICB)】God said, “Let there be light!” And there was light. ...Then God said, “Let there be something to divide the water in two!” So God made the air to divide the water in two. ...Then God said, “Let the water under the sky be gathered together so the dry land will appear.” And it happened.\\
【Hebrews 9:22 (ICB)】 The law says that almost everything must be made clean by blood. And sins cannot be forgiven without blood to show death.\\
【1 John 1:7 (ICB)】 God is light, and in him there is no darkness at all. And when we live in the light, the blood of the death of Jesus, God’s Son, is making us clean from every sin.\\
\vspace{.25cm}
\textbf{Questions}\vspace{-.5cm}
\begin{itemize}[noitemsep]
\item According to Genesis, how did God create the universe?
\item According to the Law, what should be used to make everything clean?
\item How can sins be forgiven?
\item Whose blood of the death can make us clean from every sin?
\item What is light according to John?
\item What should we do to be cleaned?
\end{itemize}

\vspace{.25cm}
\subsubsection{\large{Mark 2:1-3}}
A few days later, Jesus came back to Capernaum. The news spread that he was home. So many people gathered to hear him preach that the house was full. There was no place to stand, not even outside the door. Jesus was teaching them. Some people came, bringing a paralyzed man to Jesus. Four of them were carrying the paralyzed man. \\
\vspace{.25cm}
\textbf{Questions}\vspace{-.5cm}
\begin{itemize}[noitemsep]
\item “he was home”， here, "home" refers to which city?
\item Where did Jesus give his preach?
\item Why there was no place to stand?
\item Why the man was carried by others?
\item How many people carried the man?
\end{itemize}

\vspace{.25cm}
\subsubsection{\large{Mark 2:4-5}}
But they could not get to Jesus because of the crowd. So they went to the roof above Jesus and made a hole in the roof. Then they lowered the mat with the paralyzed man on it. Jesus saw that these men had great faith. So he said to the paralyzed man, “Young man, your sins are forgiven.” \\
\vspace{.25cm}
\textbf{Questions}\vspace{-.5cm}
\begin{itemize}[noitemsep]
\item How did the four man send the paralyzed man to Jesus?
\item Whose great faith was found by Jesus? From which description you draw the conclusion?
\item The paralyzed man was a  \rule[0cm]{1cm}{0.4pt}\\
A, old man \hspace{1cm} B, child, \hspace{1cm} C, young man
\item What did Jesus say to the paralyzed man?
\end{itemize}

\newpage
\vspace{-.3cm}
\subsubsection{\large{Mark 2:6-7}}
Some of the teachers of the law were sitting there. They saw what Jesus did, and they said to themselves, “Why does this man say things like that? He is saying things that are against God. Only God can forgive sins.” \\
\vspace{.2cm}
\textbf{Questions}\vspace{-.5cm}
\begin{itemize}[noitemsep]
\item Who were sitting with Jesus in the room, and said to themselves?
\item“Why does this man say things like that"? here, "\textbf{this man}“ refers to whom?
\item“Why does this man say things like that"? here, "\textbf{things}“ refers to what?
\item Why did they disagree with Jesus?
\item When these teachers thought "Only God can forgive sins", who did they think Jesus is?
\end{itemize}


\vspace{-.2cm}
\subsubsection{\large{Mark 2:8-12}}
At once Jesus knew what these teachers of the law were thinking. So he said to them, “Why are you thinking these things? Which is easier: to tell this paralyzed man, ‘Your sins are forgiven,’ or to tell him, ‘Stand up. Take your mat and walk’? But I will prove to you that the Son of Man has authority on earth to forgive sins.” So Jesus said to the paralyzed man, “I tell you, stand up. Take your mat and go home.” Immediately the paralyzed man stood up. He took his mat and walked out while everyone was watching him. The people were amazed and praised God. They said, “We have never seen anything like this!”\\ 
%\vspace{.25cm}
\textbf{Questions}\vspace{-.25cm}
\begin{itemize}[noitemsep]
\item Did Jesus know what's the thought in these teachers's mind?
\item "I will prove to you that the Son of Man has authority on earth", here \textbf{the Son of Man} refers to whom? The word of \textbf{authority} refers to which kind of authority on earth?
\item The paralyzed man was brought to Jesus, what was this man's epectation from Jesus? Did Jesus know this man's expectation?
\item However, what \textbf{did} Jesus want to give the paralyzed man prior to giving him the healling?
\item Why did Jesus say： “Young man, your sins are forgiven” rather than say ‘Stand up. Take your mat and walk’ first? (Use the word of Bible to answer the question)
\item Jesus asked a question to people in the house, but Jesus did not give the answer. What is your answer about Jesus's question. 【Which is easier: to tell this paralyzed man, ‘\textbf{Your sins are forgiven},’ or to tell him, ‘\textbf{Stand up. Take your mat and walk}’?】
\item If \underline{a man} tells the paralyzed man, ‘\textbf{Stand up. Take your mat and walk}’, how can others determin his statement has the power of healing or not?
\item If \underline{a man} tells the paralyzed man, ‘\textbf{your sins are forgiven}’, can others easily determin his word has the power of forgiven or not? or his statement is "true" or "false"?
\item For Jesus, what he has to do to make ‘\textbf{your sins are forgiven}’ become TRUE?
\item What is the price Jesus has to pay to get the power of "forgiven sins"?
\item How did Jesus prove to the people on site that he really has the power of  "forgiven sins"? or why did the sequence of Jesus' word could prove that he has the power of forgive sins?
\item How did God create the universe? What did Jesus do to heal the paralyzed man?
\item \textbf{For Jesus} which is easier: to tell this paralyzed man, ‘\textbf{Your sins are forgiven},’ or to tell him, ‘\textbf{Stand up. Take your mat and walk}’?
\item \textbf{For human being} which is easier: to tell this paralyzed man, ‘\textbf{Your sins are forgiven},’ or to tell him, ‘\textbf{Stand up. Take your mat and walk}’?
\item What were the people's reactions after seeing these?
\end{itemize}















\newpage
\Large
\subsection{\LARGE{Youth Version - English Version}}
\large
\vspace{.25cm}
\subsubsection{\large{Genesis 1:3,6-7,9/Hebrews 9:22/1 John 1:7 (ICB)}}
【Genesis 1:3,6-7,9 (ICB)】God said, “Let there be light!” And there was light. ...Then God said, “Let there be something to divide the water in two!” So God made the air to divide the water in two. ...Then God said, “Let the water under the sky be gathered together so the dry land will appear.” And it happened.\\
【Hebrews 9:22 (ICB)】 The law says that almost everything must be made clean by blood. And sins cannot be forgiven without blood to show death.\\
【1 John 1:7 (ICB)】 God is light, and in him there is no darkness at all. And when we live in the light, the blood of the death of Jesus, God’s Son, is making us clean from every sin.\\
\vspace{.25cm}
\textbf{Questions}\vspace{-.5cm}
\begin{itemize}[noitemsep]
\item According to Genesis, how did God create the universe?
\item According to the Law, what should be used to make everything clean?
\item How can sins be forgiven?
\item Whose blood of the death can make us clean from every sin?
\item What is light according to John?
\item What should we do to be cleaned?
\end{itemize}

\vspace{.25cm}
\subsubsection{\large{Mark 2:1-7}}
A few days later, Jesus came back to Capernaum. The news spread that he was home. So many people gathered to hear him preach that the house was full. There was no place to stand, not even outside the door. Jesus was teaching them. Some people came, bringing a paralyzed man to Jesus. Four of them were carrying the paralyzed man. But they could not get to Jesus because of the crowd. So they went to the roof above Jesus and made a hole in the roof. Then they lowered the mat with the paralyzed man on it. Jesus saw that these men had great faith. So he said to the paralyzed man, “Young man, your sins are forgiven.” Some of the teachers of the law were sitting there. They saw what Jesus did, and they said to themselves, “Why does this man say things like that? He is saying things that are against God. Only God can forgive sins.” \\
\vspace{.25cm}
\textbf{Questions}\vspace{-.5cm}
\begin{itemize}[noitemsep]
\item “he was home”， here, "home" refers to which city?
\item Where did Jesus give his preach?
\item Why there was no place to stand?
\item Why the man was carried by others? How many people carried the man?
\item How did the four man send the paralyzed man to Jesus?
\item Whose great faith was found by Jesus? From which description you draw the conclusion?
\item The paralyzed man was a  \rule[0cm]{1cm}{0.4pt}\\
A, old man \hspace{1cm} B, child, \hspace{1cm} C, young man
\item What did Jesus say to the paralyzed man?
\item Who were sitting with Jesus in the room, and said to themselves?
\item“Why does this man say things like that"? here, "\textbf{this man}“ refers to whom?
\item“Why does this man say things like that"? here, "\textbf{things}“ refers to what?
\item Why did they disagree with Jesus?
\item When these teachers thought "Only God can forgive sins", who did they think Jesus was?
\end{itemize}


\vspace{.3cm}
\subsubsection{\large{Mark 2:8-12}}
At once Jesus knew what these teachers of the law were thinking. So he said to them, “Why are you thinking these things? Which is easier: to tell this paralyzed man, ‘Your sins are forgiven,’ or to tell him, ‘Stand up. Take your mat and walk’? But I will prove to you that the Son of Man has authority on earth to forgive sins.” So Jesus said to the paralyzed man, “I tell you, stand up. Take your mat and go home.” Immediately the paralyzed man stood up. He took his mat and walked out while everyone was watching him. The people were amazed and praised God. They said, “We have never seen anything like this!”\\ 
\vspace{.25cm}
\textbf{Questions}\vspace{-.5cm}
\begin{itemize}[noitemsep]
\item Did Jesus know what's the thought in these teachers's mind?
\item "I will prove to you that the Son of Man has authority on earth", here \textbf{the Son of Man} refers to whom? The word of \textbf{authority} refers to which kind of authority on earth?
\item The paralyzed man was brought to Jesus, what was this man's epectation from Jesus? Did Jesus know this man's expectation?
\item However, what \textbf{did} Jesus want to give the paralyzed man prior to giving him the healling?
\item Why did Jesus say： “Young man, your sins are forgiven” rather than say ‘Stand up. Take your mat and walk’ first? (Use the word of Bible to answer the question)
\item Jesus asked a question to people in the house, but Jesus did not give the answer. What is your answer about Jesus's question. 【Which is easier: to tell this paralyzed man, ‘\textbf{Your sins are forgiven},’ or to tell him, ‘\textbf{Stand up. Take your mat and walk}’?】
\item If \underline{a man} tells the paralyzed man, ‘\textbf{Stand up. Take your mat and walk}’, how can others determin his statement has the power of healing or not?
\item If \underline{a man} tells the paralyzed man, ‘\textbf{your sins are forgiven}’, can others easily determin his word has the power of forgiven or not? or his statement is "true" or "false"?
\item For Jesus, what he has to do to make ‘\textbf{your sins are forgiven}’ become TRUE?
\item What is the price Jesus has to pay to get the power of "forgiven sins"?
\item How did Jesus prove to the people on site that he really has the power of  "forgiven sins"?
\item How did God create the universe? What did Jesus do to heal the paralyzed man?
\item \textbf{For Jesus} which is easier: to tell this paralyzed man, ‘\textbf{Your sins are forgiven},’ or to tell him, ‘\textbf{Stand up. Take your mat and walk}’?\\
\textbf{For human being} which is easier: to tell this paralyzed man, ‘\textbf{Your sins are forgiven},’ or to tell him, ‘\textbf{Stand up. Take your mat and walk}’?
\item What were the people's reactions after seeing these?
\end{itemize}














\iffalse
And again Jesus entered Capernaum, and it was heard that He was in the house. 2 Immediately many gathered together, so that there was no longer room to receive them, not even near the door. And He preached the word to them. 3 Then they came to Him, bringing a paralytic who was carried by four men. 4 And when they could not come near Him because of the crowd, they uncovered the roof where He was. So when they had broken through, they let down the bed on which the paralytic was lying.

5 When Jesus saw their faith, He said to the paralytic, “Son, your sins are forgiven you.”

6 And some of the scribes were sitting there and reasoning in their hearts, 7 “Why does this Man speak blasphemies like this? Who can forgive sins but God alone?”

8 But immediately, when Jesus perceived in His spirit that they reasoned thus within themselves, He said to them, “Why do you reason about these things in your hearts? 9 Which is easier, to say to the paralytic, ‘Your sins are forgiven you,’ or to say, ‘Arise, take up your bed and walk’? 10 But that you may know that the Son of Man has [b]power on earth to forgive sins”—He said to the paralytic, 11 “I say to you, arise, take up your bed, and go to your house.” 12 Immediately he arose, took up the bed, and went out in the presence of them all, so that all were amazed and glorified God, saying, “We never saw anything like this!”







And when he entered again into Capernaum after some days, it was noised that he was in the house. 2And many were gathered together, so that there was no longer room for them, no, not even about the door: and he spake the word unto them. 3And they come, bringing unto him a man sick of the palsy, borne of four. 4And when they could not come nigh unto him for the crowd, they uncovered the roof where he was: and when they had broken it up, they let down the bed whereon the sick of the palsy lay. 5And Jesus seeing their faith saith unto the sick of the palsy, Son, thy sins are forgiven. 6But there were certain of the scribes sitting there, and reasoning in their hearts, 7Why doth this man thus speak? he blasphemeth: who can forgive sins but one, even God? 8And straightway Jesus, perceiving in his spirit that they so reasoned within themselves, saith unto them, Why reason ye these things in your hearts? 9Whether is easier, to say to the sick of the palsy, Thy sins are forgiven; or to say, Arise, and take up thy bed, and walk? 10But that ye may know that the Son of man hath power on earth to forgive sins (he saith to the sick of the palsy), 11I say unto thee, Arise, take up thy bed, and go unto thy house. 12And he arose, and straightway took up the bed, and went forth before them all; insomuch that they were all amazed, and glorified God, saying, We never saw it on this fashion.
\fi